% lem-radial -- anchor CERW.Frozen.radial_test. Source: coarse-radius.tex:38-58 (flat 581-601); b_d, M_sh, phi_r at 20-36.
\paragraph{Paper (verbatim).} ``Let $b_d=\lceil\sqrt d/2\rceil+6$, and let
$F(s)=\sigma_d^{-1}\int_{D_n\cap\{|v|>s\}}|v|^{1-d}\dd v$. For every $p>0$, there are $r_0=r_0(d)>2b_d$
and $C=C(d,\eps,p)>0$ such that, for every integer $n\geq2$, with probability at least $1-Cn^{-p}$,
simultaneously for all integers $r_0\leq r\leq n$,
$F(r+b_d)\leq CM_{\rm sh}(r)[F(r-b_d)-F(r+b_d)]+C[\sqrt{M_nr^{1-d}F(r-b_d)L}+r^{1-d}L]$.''
\paragraph{Proof.} Let $h_d(r)=\frac2\pi\log r+\kappa$ ($d=2$) or $-\frac2{(d-2)\omega_d}r^{2-d}$
($d\ge3$), and $\phi_r=(b-h_d(r))_+$ with $b$ from \texttt{s-kernel-props}.
\emph{Step 1 (level sets).} By \texttt{s-levelsets} there is $r_0=r_0(d)>2b_d$, depending only on the
External's constants, such that for $r\ge r_0$: $\phi_r=0$ on $|x|\le r-1$ and $\phi_r=b-h_d(r)$ on
$|x|\ge r+1$. The radial profile gains order $r^{1-d}$ over a unit step, against an $O(r^{-d})$
error. Inner sites are covered by $h_2\to\infty$, and by $h_d\to0$ with $b<0$ when $d\ge3$.
\emph{Step 2 (drift).} \texttt{s-radial-drift}: the increments of $\phi_r$ vanish on $|x|\le r-2$.
For $|x|>r+3$, $(P-I)\phi_r(x)=0$ and $u_x\cdot D\phi_r(x)\ge c_d|x|^{1-d}$ (\texttt{s-gradient}).
On the shell $||x|-r|\le3$, $|(P-I)\phi_r|,|D\phi_r|\le C_dr^{1-d}$, because the positive part is
1-Lipschitz.
\emph{Step 3 (martingale).} \texttt{s-dynkin} with $f=\phi_r$ gives the martingale $W^{(r)}$, with
increments $\le C_dr^{1-d}$ and bracket $\le C_d\sum_{|x|>r-2}\ell_n(x)|x|^{2-2d}\le C_dM_nr^{1-d}F(r-b_d)$.
Each such departure cell lies outside $r-b_d$, and its weighted integral is comparable to
$|x|^{1-d}$ (\texttt{s-Fmass}, \texttt{s-cell-norm}). \texttt{s-dyadic}, applied to $r^{d-1}W^{(r)}$
over the $n$ integer radii and the dyadic levels (including the level below one), gives an event
$G$ with $\mu(G^c)\le Cn^{-p}$ on which
$|W^{(r)}_n|\le C[\sqrt{M_nr^{1-d}F(r-b_d)L}+r^{1-d}L]$ for all $r_0\le r\le n$
(\texttt{s-radial-mart}).
\emph{Step 4 (source).} Since $\phi_r(X_0)=\phi_r(0)=0$ and $\phi_r(X_n)\ge0$, Dynkin gives
$c_d\eps\sum_{x\in A_n,|x|>r+3}|x|^{1-d}\le C_dr^{1-d}\sum_{||x|-r|\le3}\ell_n(x)+|W^{(r)}_n|$. The
first-departure shell term is absorbed by the local time, since each shell site of $A_n$ has
$\ell_n\ge1$. The left side is at least $cF(r+b_d)$, because each cell meeting $\{|v|>r+b_d\}$ has
$|x|>r+3$. The shell sum is at most $M_{\rm sh}(r)$ times the number of shell sites of $A_n$, which
is at most $\sigma_d(r+b_d)^{d-1}[F(r-b_d)-F(r+b_d)]$: shell cells lie in $r-b_d<|v|\le r+b_d$
(\texttt{s-Fmass}, increments). Multiplying through by $r^{1-d}$ gives the claim on $G$
(\texttt{s-radial-source}).
\paragraph{Transcription notes.} $r$ ranges over naturals $r_0\le r\le n$, with $r_0>2b_d$. The
arguments $r\pm b_d$ are real ($(r:\R)\pm b_d$). $M_{\rm sh}(r)$ is \texttt{shellMax}, a finite
\texttt{Finset.sup}. $r_0$ is bound before $\eps$ (it depends only on $d$), and $C$ after $\eps,p$.
